% !TeX program = XeLaTeX
% !TeX root = ../vedamantrabook.tex


\chapt{भाग्यसूक्तम्}
\vspace{-1ex}
\centerline{\scriptsize(तैत्तिरीय ब्राह्मणम् अष्टकम् – ३/प्रश्नः – ८/अनुवाकम् – ९)}

\begingroup\newcommand{\ShlokaStrut}{\vphantom{र्क॑हु॒}}

\twolineshloka
{\ShlokaStrut{}प्रा॒तर॒ग्निं प्रा॒तरिन्द्रꣳ॑ हवामहे प्रा॒तर्मि॒त्रा वरु॑णा प्रा॒तर॒श्विना᳚}
{\ShlokaStrut{}प्रा॒तर्भगं॑ पू॒षणं॒ ब्रह्म॑ण॒स्पतिं॑ प्रा॒तः सोम॑मु॒त रु॒द्रꣳ हु॑वेम}

\twolineshloka
{\ShlokaStrut{}प्रा॒त॒र्जितं॒ भग॑मु॒ग्रꣳ हु॑वेम व॒यं पु॒त्रमदि॑ते॒र्यो वि॑ध॒र्ता}
{\ShlokaStrut{}आ॒र्ध्रश्चि॒द्यं मन्य॑मानस्तु॒रश्चि॒द्राजा॑ चि॒द्यं भगं॑ भ॒क्षीत्याह॑}

\twolineshloka
{\ShlokaStrut{}भग॒ प्रणे॑त॒र्भग॒ सत्य॑राधो॒ भगे॒मां धिय॒मुद॑व॒ दद॑न्नः}
{\ShlokaStrut{}भग॒ प्र णो॑ जनय॒ गोभि॒रश्वै॒र्भग॒ प्र नृभि॑र्नृ॒वन्तः॑ स्याम}

\twolineshloka
{\ShlokaStrut{}उ॒तेदानीं॒ भग॑वन्तः स्यामो॒त प्रपि॒त्व उ॒त मध्ये॒ अह्ना᳚म्}
{\ShlokaStrut{}उ॒तोदि॑ता मघव॒न्थ्सूर्य॑स्य व॒यं दे॒वानाꣳ॑ सुम॒तौ स्या॑म}

\twolineshloka
{\ShlokaStrut{}भग॑ ए॒व भग॑वाꣳ अस्तु देवा॒स्तेन॑ व॒यं भग॑वन्तः स्याम}
{\ShlokaStrut{}तं त्वा॑ भग॒ सर्व॒ इज्जो॑हवीमि॒ स नो॑ भग पुर ए॒ता भ॑वे॒ह}

\twolineshloka
{\ShlokaStrut{}सम॑ध्व॒रायो॒षसो॑ऽनमन्त दधि॒क्रावे॑व॒ शुच॑ये प॒दाय॑}
{\ShlokaStrut{}अ॒र्वा॒ची॒नं व॑सु॒विदं॒ भगं॑ नो॒ रथ॑मि॒वाश्वा॑वा॒जिन॒ आव॑हन्तु}

\twolineshloka
{\ShlokaStrut{}अश्वा॑वती॒र्गोम॑तीर्न उ॒षासो॑ वी॒रव॑तीः॒ सद॑मुच्छन्तु भ॒द्राः}
{\ShlokaStrut{}घृ॒तं दुहा॑ना वि॒श्वतः॒ प्रपी॑ना यू॒यं पा॑त स्व॒स्तिभिः॒ सदा॑ नः}

\twolineshloka
{\ShlokaStrut{}यो मा᳚ऽग्ने भा॒गिनꣳ॑ स॒न्तमथा॑भा॒गं चिकी॑र्‌षति}
{\ShlokaStrut{}अ॒भा॒गम॑ग्ने॒ तं कु॑रु॒ माम॑ग्ने भा॒गिनं॑ कुरु}

\endgroup